\chapter{Discussion}
\label{chapter:discussion}

This chapter interprets and contextualizes the experimental results presented in Chapter~\ref{chapter:results}. Rather than reiterating numerical findings, the discussion focuses on explaining observed system behavior, analyzing trade-offs inherent to real-time simultaneous interpretation, and assessing the robustness of the proposed architecture under different configurations and parameter settings. The discussion is grounded in the research objectives defined in Chapter~\ref{chapter:introduction} and follows the evaluation framework outlined in Chapter~\ref{chapter:methodology}. Moreover, it is structured to mirror the logical progression of the results. It first examines latency characteristics of the cascaded interpretation pipeline at both end-to-end and component levels, before analyzing the stability--latency trade-offs revealed through ablation studies. Finally, the limitations of the evaluation are discussed to clearly delineate the scope and applicability of the findings.

\section{Latency Characteristics of the Cascaded Interpretation Architecture}

A central objective of this thesis was to design a real-time interpretation system capable of maintaining low and stable latency under realistic operating conditions. The end-to-end latency metrics reported in Section~\ref{sec:e2e_latency} demonstrate that the proposed system consistently operates within latency bounds commonly considered acceptable for simultaneous interpretation \cite{ ma2019staclarxiv, fantinuoli2023defining}. Across all evaluated configurations, AL, LAAL, and RTF remained stable with only minor variations between baseline and alternative component configurations. The absence of large latency fluctuations across configurations is a significant result. While swapping ASR or TTS components introduces local differences in processing behavior, these changes do not propagate uncontrollably through the pipeline. Instead, the orchestration logic and segmentation mechanisms absorb component-level variability, resulting in consistent end-to-end behavior. From an architectural perspective, this indicates that the system design effectively decouples individual processing stages while preserving global temporal coherence. This robustness is particularly relevant in practical deployment scenarios, where AI service providers may be changed due to cost, availability, or language coverage considerations. The results suggest that such substitutions can be done without compromising simultaneity, provided that the components conform to the expected streaming interfaces and latency profiles.

%\subsection{Distribution of Latency Across Pipeline Stages}

Component-level latency analysis provides insight into how end-to-end delay emerges from the interaction of individual pipeline stages. As reported in Section~\ref{sec:component_level_latency}, the ASR and segmentation stage consistently accounts for the largest share of total latency across all evaluated configurations. This result is expected, as streaming recognition operates on continuously evolving acoustic input and must reconcile low-latency emission with the need to suppress unstable or incomplete hypotheses. The segmentation mechanism further contributes to this delay by intentionally postponing output finalization until stability criteria are met, thereby trading immediacy for improved downstream consistency.

In contrast, the NMT stage contributes a smaller and comparatively stable portion of the overall latency budget. This behavior reflects the relatively predictable inference characteristics of modern neural machine translation systems when operating on short, incrementally generated text segments. Translation latency remains well bounded across baseline and alternative configurations, indicating that NMT inference does not constitute a dominant bottleneck under the tested conditions and is less sensitive to upstream segmentation variability.

The TTS stage introduces a substantial but controlled delay associated with speech synthesis and audio playback initiation. While neural synthesis is computationally demanding, its latency remains consistent due to the use of fixed buffering and scheduling strategies designed for real-time output. Notably, substituting alternative TTS components does not lead to disproportionate increases in synthesis latency, reinforcing the observation that the system architecture effectively constrains component-level variability and preserves predictable temporal behavior at the output stage.

%\subsection{Implications for System Optimization}

The observed distribution of latency across pipeline stages has direct implications for system optimization. Since ASR and segmentation dominate the latency budget, improvements in recognition responsiveness, hypothesis stabilization, or segmentation policy design are likely to yield the most significant reductions in end-to-end delay. In comparison, further optimization of the NMT stage would provide limited benefit, given its relatively small and stable contribution to overall latency.

These findings emphasize the importance of adopting a system-level optimization perspective rather than focusing exclusively on individual model performance. Effective orchestration, buffering strategies, and segmentation policies play a critical role in shaping real-time behavior and should be treated as first-class design elements alongside the selection of ASR, NMT, and TTS components. This perspective is particularly relevant for practical deployment scenarios, where predictable latency and robustness are often more important than marginal gains in isolated component efficiency.


\section{Stability--Latency Trade-Offs in Incremental Segmentation}

Beyond absolute latency, perceived quality in simultaneous interpretation is strongly influenced by output stability. The ablation study on stability window duration exposes a clear trade-off between responsiveness and output smoothness. As expected, increasing the fixed stability window consistently reduces flicker rate and normalized erasure, as segment finalization is delayed until recognition hypotheses converge and become less susceptible to revision.

A notable and unexpected observation is that the adaptive stability baseline does not achieve the lowest instability metrics. Despite dynamically adjusting the stability window based on speaker rate, the adaptive strategy exhibits higher flicker and erasure than some fixed-window configurations. This behavior can be explained by local oscillations in the adaptation mechanism, where frequent adjustments to the stability window introduce additional revisions at short time scales, thereby increasing provisional output volatility.

When evaluated at the session level, however, the adaptive strategy yields the lowest mean end-to-end latency. This highlights an important distinction between local stability measures and global system behavior. While adaptive segmentation may introduce short-term instability, it avoids overly conservative buffering during fluent speech, thereby improving overall responsiveness. The results demonstrate that minimizing flicker or erasure in isolation does not necessarily correspond to optimal simultaneity or perceived performance in live interpretation scenarios. Therefore, from a deployment perspective, these findings suggest that stability parameters should be tuned with respect to session-level behavior rather than individual instability metrics alone. A moderate degree of provisional output revision may be acceptable if it leads to more consistent temporal alignment between source and target speech. This reinforces the need for holistic evaluation of segmentation strategies that jointly considers latency, stability, and listener experience, rather than optimizing any single metric in isolation.


\section{Effects of Segmentation Aggressiveness and Token Thresholding}

The ablation study on dynamic token thresholds further illustrates the non-linear behavior of segmentation mechanisms in streaming interpretation systems. Contrary to initial expectations, lower token thresholds do not consistently reduce CW. Instead, highly aggressive segmentation increases buffering frequency and coordination overhead between pipeline stages, resulting in higher CW values and elevated segment-level latency.

Frequent forced emissions introduce additional synchronization points across recognition, translation, and synthesis components. Rather than improving responsiveness, this behavior amplifies processing overhead and destabilizes timing alignment. In extreme cases, aggressive token thresholds lead to segment-level delays that exceed acceptable bounds for simultaneous interpretation, despite their intended role in minimizing latency. These results indicate that emission frequency alone is an insufficient proxy for responsiveness in cascaded real-time systems. Therefore, from a system perspective, the results suggest that more conservative token thresholding yields more stable and predictable behavior. While aggressive segmentation appears advantageous in theory, its interaction with stabilization logic, buffering mechanisms, and downstream processing can degrade both temporal coherence and overall system efficiency. This underscores the importance of evaluating segmentation aggressiveness within the context of the complete pipeline, rather than optimizing token-level policies in isolation.

\section{Limitations of the Evaluation and Scope Considerations}

While the presented results provide meaningful insight into the behavior of the proposed system, several limitations must be acknowledged to contextualize the evaluation. First, the experiments are based on scripted speech produced by a single speaker. This setup reduces variability in speaking style, accent, prosody, and disfluencies. Although this choice enables controlled, repeatable measurements and facilitates fair comparison across configurations, it does not fully reflect the diversity of spontaneous speech encountered in real-world interpretation scenarios, where speaker behavior may vary significantly.

Second, the evaluation deliberately focuses on system-level performance metrics, such as latency, stability, and segmentation behavior, rather than on linguistic translation quality. Measures of semantic accuracy, fluency, or adequacy are therefore not included. This reflects the primary objective of the thesis, which is to investigate architectural feasibility, real-time orchestration, and latency control in AI-based simultaneous interpretation systems, rather than to benchmark or optimize individual language models.

Finally, all experiments were conducted under controlled hardware and network conditions. Stable network connectivity, consistent audio capture, and fixed computational resources were assumed to isolate system behavior from external disturbances. As a result, the reported performance may differ in less predictable environments, such as wireless networks, mobile deployments, or large-scale distributed setups.

These limitations represent deliberate scope decisions rather than methodological shortcomings. By constraining experimental conditions, the evaluation isolates the impact of architectural and segmentation choices and ensures interpretability of the results. The findings therefore provide a solid foundation for future work, which could extend the evaluation to multi-speaker scenarios, spontaneous speech, and qualitative user-centered assessments.

\noindent
This thesis has addressed the research objectives defined at the outset by demonstrating the feasibility of a real-time AI-based simultaneous interpretation system from a system-engineering perspective. The proposed modular pipeline successfully integrates streaming ASR, NMT, and TTS synthesis while maintaining end-to-end latency within bounds suitable for live interpretation scenarios. Across all evaluated configurations, the system exhibits stable latency behavior, indicating that the architectural design effectively absorbs variability introduced by alternative ASR and TTS components. This robustness supports the central premise that careful orchestration and segmentation strategies are at least as critical as the choice of individual AI models for achieving reliable real-time performance.

The component-level latency analysis further clarifies how delays accumulate across the pipeline and highlights the dominant role of recognition and segmentation in shaping overall system behavior. While text-to-speech synthesis contributes a substantial share of latency, its impact remains bounded due to controlled buffering and scheduling mechanisms. Machine translation latency, by contrast, remains comparatively stable across configurations. These observations reinforce the importance of system-level optimization, particularly in upstream stages where segmentation and stabilization decisions directly influence both responsiveness and output smoothness. The results suggest that future improvements in real-time interpretation systems are more likely to arise from refined orchestration and segmentation policies than from isolated model-level optimization.

Finally, the ablation studies reveal nuanced trade-offs between stability, latency, and segmentation aggressiveness. Increasing the stability window reduces output instability, as expected, yet the adaptive stability strategy does not consistently minimize flicker or erasure. Despite this, it achieves the lowest mean end-to-end latency over extended sessions, illustrating that local instability metrics do not fully capture perceived system performance. Similarly, aggressive token-based segmentation does not necessarily improve responsiveness; instead, frequent emissions can increase buffering and coordination overhead, leading to higher effective latency. These findings underscore the complexity of incremental processing and demonstrate that segmentation parameters must be tuned holistically rather than optimized in isolation. Taken together, the results highlight both the promise and the challenges of AI-based simultaneous interpretation and provide a grounded basis for future work exploring broader deployment scenarios, richer evaluation protocols, and extended user-centered assessments.